\section{Discussion}\label{sec:discussion}

Separating surface mass balance from dynamic thinning in an observed
surface-elevation change is a notoriously ill-posed problem, yet it is exactly what
deriving SMB from remote sensing demands. We have addressed it with a transient,
differentiable inversion: a state-of-the-art assimilation of surface ice speed fixes
the dynamical state, and a novel automatic-differentiation scheme then inverts the
multi-temporal geometry change for the elevation-resolved SMB. On three Alpine
glaciers of contrasting size the method reproduces the in situ stake balances to
root-mean-square errors of $0.40$ (Argentière), $0.52$ (Great Aletsch) and
$0.60\,\mathrm{m\,w.e.\,yr^{-1}}$ (Rhône), surpassing the best physically based
reconstruction of \citet{kneib2024} on the same Argentière stakes with surface
velocity alone. Two features set the approach apart from that state of the art.
First, the flux divergence is recomputed from the evolving geometry at every time step
rather than held at a single stationary snapshot. A stationary estimate evaluates
$\nabla\!\cdot\mathbf{q}$ once, on the geometry of a single epoch, and applies it to an
elevation change averaged over a decade---a decade over which the glacier thins, slows
and changes its flux. Integrating the mass-conservation equation forward removes that
assumption rather than correcting for it, so the inferred SMB is dynamically consistent
across the whole window; against the same stakes, the stationary-flux estimate is
markedly worse on all three glaciers (Sect.~\ref{sec:results}).

Second, the inversion needs only surface velocity---no
ice-thickness survey and no per-glacier tuning---and proved largely insensitive to
the rate factor $A$, so a single glacier-independent pipeline transfers across all
three sites. Beyond validating the inversion, the recovered $smb(z)$ is the quantity
mass-balance models are calibrated against: an elevation-resolved balance over a known
period is what a temperature-index or energy-balance model must reproduce before it can
be driven forward under a climate scenario. Inferring it from geometry alone therefore
offers a route to constraining such models---and with them, projections of glacier
change and meltwater supply---on the great majority of glaciers that carry no stake
record.

The clearest methodological gain is in how mass conservation is enforced.
Flux-divergence inversions form the SMB by differentiating the observed
thickness--velocity flux and combining it with the elevation-change field; because
that divergence is the spatial derivative of a product of two noisy fields, it must
be smoothed before use, and the smoothing---as \citet{miles2021} and \citet{kneib2024}
document---breaks
the very mass-conservation constraint the method rests on, leaving a spurious
glacier-wide flux-divergence residual of order $0.1$--$0.7\,\mathrm{m\,yr^{-1}}$ that
they close by redistributing the excess or deficit uniformly across the glacier. Such
a homogeneous redistribution is a pragmatic fix rather than a physical statement: it
spreads a computational residual over cells that did not produce it. Our formulation
removes the need for either step. Because we integrate the mass-conservation equation
forward and evaluate the flux divergence on the modelled geometry with a
conservative, slope-limited scheme, mass is conserved by construction at every step;
the inferred SMB field requires no post-hoc smoothing and no ad-hoc redistribution,
and the divergence it implies is guaranteed consistent with the simulated geometry
change.

The contrast is visible in our own runs. The flux divergence implied by the calibrated
assimilation state is markedly rougher than the thickness field it derives from---
grid-scale structure carries three to four times the share of its variance---because
differentiating a product of two fields amplifies exactly the wavelengths the
assimilation constrains least. Integrating forward damps it: the flux is diffusive in
thickness, so by the end of the window the modelled divergence is smooth and spatially
coherent where the assimilated one was speckled. The filtering that flux-divergence
methods must impose externally is here a consequence of the ice dynamics rather than a
step applied afterwards.

The method's principal limitation is its one-dimensional parametrisation: horizontal
variability within an elevation band---such as the avalanche-fed accumulation
hotspots resolved by \citet{kneib2024}---is deliberately not recovered. This buys a
transferable, tuning-free inversion and is not intrinsic to the scheme, which admits
richer descriptions (an aspect or debris-cover dependence, say) wherever the data
warrant them. A second limitation is the thickness--sliding ambiguity of
Sect.~\ref{sec:sensitivity}: surface velocity constrains the basal-motion parameter
only to a band, and the resulting freedom, negligible in the accumulation area, opens
into an equilibrium-line spread at the tongue. Two further dependencies are shared
with flux-based methods. Accuracy is ultimately set by the input velocity and DEMs:
\citet{kneib2024} find that substituting the coarser global velocity mosaic of
\citet{millan2022} into their reconstruction degrades the stake agreement to
$\sim$1.2--1.3$\,\mathrm{m\,w.e.\,yr^{-1}}$, so the velocity field can itself become
limiting---that our transient inversion attains $0.40$--$0.60\,\mathrm{m\,w.e.\,yr^{-1}}$
using precisely that global product is encouraging, but its robustness to velocity
error should not be overstated. Finally, like all geodetic estimates the method
converts a volume change to a mass balance under an assumed ice density and neglects
evolving firn compaction; where the firn column is thick or densifying, this
introduces a systematic offset.

Taken together, these properties point towards regional-scale application. Because the
pipeline requires no thickness survey, no stake calibration and no site-specific
tuning, it can in principle be applied wherever multi-temporal DEMs and a
surface-velocity field exist---and both are approaching near-global coverage, through
the elevation-change archive of \citet{hugonnet2021} and expanding velocity products
\citep{millan2022}. Its accuracy will track that of these inputs: substituting that
$\mathrm{d}h/\mathrm{d}t$ product for the dedicated DEM pair on
Rhône raises the stake error to $0.77\,\mathrm{m\,w.e.\,yr^{-1}}$ (from $0.60$),
while still
recovering the full altitudinal structure and remaining far better than the
stationary-flux estimate (Appendix~\ref{app:hugonnet}). The highest fidelity therefore
still comes from dedicated high-resolution DEMs, but the method degrades gracefully
rather than failing where only near-global data exist; and the tuning-free character
makes a continuous, distributed SMB inversion over many glaciers a realistic next
step. Whether the robustness shown here on three well-monitored Alpine glaciers
persists across the smaller and more varied population that dominates regional water
resources is the key open question, and the natural direction for future work.
